\documentclass[a4paper,10pt]{ltjsarticle}
\newcommand{\bookletname}{問題冊子}
% LuaLaTeX / UTF-8  第2回模試 共通スタイル（統計応用（医薬生物学））
\usepackage{amsmath,amssymb,bm}
\usepackage[margin=20mm,top=20mm,bottom=19mm]{geometry}
\usepackage{enumitem,booktabs,array,multirow}
\usepackage{fancyhdr}
\usepackage[hidelinks,unicode]{hyperref}
\setlength{\parindent}{1\zw}
\setlength{\parskip}{0.25em}
\setlength{\headheight}{17pt}
\pagestyle{fancy}
\fancyhf{}
\fancyhead[L]{\small 第2回模試\quad 統計応用（医薬生物学）}
\fancyhead[R]{\small\bookletname}
\fancyfoot[C]{\small\thepage}
\renewcommand{\headrulewidth}{0.35pt}
\renewcommand{\footrulewidth}{0pt}
\setlist[enumerate,1]{label=〔\arabic*〕,leftmargin=3.5em,labelsep=0.5em,itemsep=0.55em,topsep=0.5em}
\newcommand{\E}{\operatorname{E}}
\newcommand{\Var}{\operatorname{Var}}
\newcommand{\Cov}{\operatorname{Cov}}
\newcommand{\Corr}{\operatorname{Corr}}
\newcommand{\Prob}{\operatorname{P}}
\newcommand{\ind}{\mathbf{1}}
\newcommand{\qhead}[1]{\section*{問#1}}
\newcommand{\ans}[1]{\subsection*{〔#1〕}}
\newcommand{\points}[1]{\hfill\textnormal{（#1点）}}
\newcommand{\continued}[1]{\noindent\textbf{問#1（続き）}\par\medskip}
\hypersetup{pdfauthor={Claude Fable 5.2},pdfsubject={統計検定1級対策・自主学習用模擬試験}}

\hypersetup{pdftitle={第2回模試 統計応用（医薬生物学） 問題冊子}}
\setlist[enumerate,1]{label=〔\arabic*〕,leftmargin=3.5em,labelsep=0.5em,itemsep=0.3em,topsep=0.35em}
\begin{document}
\setlength{\abovedisplayskip}{5pt plus 1pt minus 2pt}\setlength{\belowdisplayskip}{5pt plus 1pt minus 2pt}\setlength{\abovedisplayshortskip}{2pt plus 1pt}\setlength{\belowdisplayshortskip}{4pt plus 1pt minus 1pt}
\thispagestyle{empty}
\begin{center}
\vspace*{12mm}
{\large 2026年 統計検定1級対策}\par\vspace{8mm}
{\Huge\bfseries 第2回 模擬試験}\par\vspace{7mm}
{\LARGE 統計応用（医薬生物学）}\par\vspace{6mm}
{\Large 問題冊子}\par\vspace{12mm}
{\large 試験時間 90分\qquad 5問中3問選択}
\end{center}
\vspace{8mm}
\noindent\textbf{受験上の注意}
\begin{enumerate}[label=\arabic*.,leftmargin=2em,itemsep=0.8em]
\item 問1から問5のうち3問を選び，選択した問題番号を答案に明記すること。解答は別紙に記述すること。
\item 答えのみでよいと明記した箇所を除き，導出または理由を示すこと。各問の記号はその問題の中だけで用いる。
\item 四則演算・平方根が計算できる一般電卓と，別添の「統計検定1級\_付表数値表」（付表1〜5）を使用してよい。本冊子には付表を再録していない。
\item $\log$ は自然対数を表す。数値は特に指定がなければ小数第3位程度まで求め，付表の丸めや補間による差は許容する。
\item 本模試では各大問を30点とする。選択した3問の合計90点を100点に換算する。小問の配点は採点基準冊子に記載する。
\item 問5は統計応用の共通問題を想定した問題である。データ・研究設定はすべて本模試のための架空のものである。
\end{enumerate}
\vfill
\begin{center}
氏名\ \underline{\hspace{55mm}}\qquad 実施日\ \underline{\hspace{35mm}}\par\vspace{8mm}
\small 自主学習用のオリジナル模試です。公式の問題・配点ではありません。\par
作成日：2026年11月12日\qquad 作成識別：fable5.2
\end{center}
\newpage

\qhead{1}
ある疾患の患者10人に新しい治療を行い，治療開始から死亡までの時間（月）を観察した。
観察された時間を小さい順に並べると
\[
2,\quad 3,\quad 4+,\quad 5+,\quad 6+,\quad 7,\quad 9,\quad 10+,\quad 12,\quad 15+
\]
であった。「$+$」は右側打ち切りを表し，打ち切りは死亡までの時間と独立に生じるものとする。

死亡が観察された時点を $t_1<t_2<\cdots<t_5$ とし，$t_j$ の直前のリスク集合の人数を $n_j$，$t_j$ での死亡数を $d_j$ とする。
生存関数 $S(t)$ のカプラン・マイヤー推定量を
\[
\hat S(t)=\prod_{j:\,t_j\le t}\left(1-\frac{d_j}{n_j}\right)
\]
とし，$G_j=\displaystyle\sum_{i=1}^{j}\frac{d_i}{n_i(n_i-d_i)}$ とおく。表1はこれらをまとめたものである。
\begin{center}
表1：カプラン・マイヤー推定の計算表\par\smallskip
\begin{tabular}{cccccc}
\toprule
$j$ & $t_j$ & $n_j$ & $d_j$ & $\hat S(t_j)$ & $G_j$\\
\midrule
1 & 2 & 10 & 1 & 0.90 & 0.0111\\
2 & 3 & （ア） & 1 & （イ） & （エ）\\
3 & 7 & 5 & 1 & 0.64 & 0.0750\\
4 & 9 & 4 & 1 & （ウ） & （オ）\\
5 & 12 & 2 & 1 & 0.24 & 0.6583\\
\bottomrule
\end{tabular}
\end{center}
以下の各問に答えよ。
\begin{enumerate}
\item 表1の（ア）〜（ウ）に入る値を求めよ。また，$\hat S(t)\le0.5$ となる最小の $t$ として，生存時間の中央値の推定値を求めよ。
\item 各 $t_j$ において，$n_j$ を与えたとき $d_j$ は二項分布 $\mathrm{Bin}(n_j,q_j)$ に従うとし，
$\hat q_j=d_j/n_j$（$j=1,2,\ldots$）は互いに無相関とみなせるとする。このとき $S(t)=\prod_{j:\,t_j\le t}(1-q_j)$ である。
デルタ法を用いて $\log\hat S(t)$ の分散の近似式を導き，$q_j$ を $\hat q_j$ で置き換えることにより，$\hat S(t)$ の分散の推定量が
\[
\widehat{\Var}\bigl[\hat S(t)\bigr]=\hat S(t)^2\sum_{j:\,t_j\le t}\frac{d_j}{n_j(n_j-d_j)}
\]
で与えられることを示せ。
\item 表1の（エ），（オ）に入る値を求めよ。また，上問〔2〕の式を用いて $\hat S(3)$ と $\hat S(9)$ の標準誤差を求め，
正規近似による $S(3)$ と $S(9)$ の95\%信頼区間をそれぞれ求めよ。
\item 時点 $t$ までに打ち切りが1件もない場合，上問〔2〕の分散の推定量は，$n$ を全体の人数として
\[
\frac{\hat S(t)\{1-\hat S(t)\}}{n}
\]
に一致することを示せ。ただし $\hat S(t)>0$ とする。
\item $S(t)$ そのものではなく $\log\{-\log S(t)\}$ について正規近似による信頼区間を作り，
それを $S(t)$ の信頼区間に変換する方法を考える。
デルタ法により $\log\{-\log\hat S(t)\}$ の分散の推定量を求め，この方法による $S(3)$ の95\%信頼区間を求めよ。
また，上問〔3〕で求めた $S(3)$ の信頼区間と比較して，この方法が優れている点を述べよ。
\end{enumerate}
\newpage

\qhead{2}
ある疾患の患者200人を治療群と対照群に100人ずつ無作為に割り付け，12週後の状態を「改善」「不変」「悪化」の3カテゴリで評価した。結果を表1に示す。
\begin{center}
表1：群別の12週後の状態（人数）\par\smallskip
\begin{tabular}{lrrrr}
\toprule
 & 改善 & 不変 & 悪化 & 計\\
\midrule
治療群 & 50 & 30 & 20 & 100\\
対照群 & 30 & 30 & 40 & 100\\
\midrule
計 & 80 & 60 & 60 & 200\\
\bottomrule
\end{tabular}
\end{center}
一般に，第 $i$ 群（$i=1,2$）の人数を $m_i$，第 $i$ 群でカテゴリ $j$（$j=1,2,3$）に属する人数を $n_{ij}$ とし，
$c_j=n_{1j}+n_{2j}$，$N=m_1+m_2$，$e_{ij}=m_ic_j/N$ とおく。以下の各問に答えよ。
\begin{enumerate}
\item 表1について，2群で3カテゴリの分布が等しいという帰無仮説 $H_0$ を，ピアソンのカイ二乗統計量を用いて有意水準5\%で検定せよ。
\item $\hat p_{ij}=n_{ij}/m_i$ とおく。次の等式を示せ。
\[
n_{1j}-e_{1j}=\frac{m_1m_2}{N}\,(\hat p_{1j}-\hat p_{2j})
\]
また，$n_{1j}$ と $n_{2j}$ が独立にそれぞれ二項分布 $\mathrm{Bin}(m_1,\pi_j)$，$\mathrm{Bin}(m_2,\pi_j)$ に従うとき（$H_0$ の下での分布），$n_{1j}-e_{1j}$ の分散を求めよ。
さらに，この分散に含まれる $\pi_j$ を $c_j/N$ で置き換えて $n_{1j}-e_{1j}$ を標準化した量
\[
r_{1j}=\frac{n_{1j}-e_{1j}}{\sqrt{e_{1j}\left(1-\dfrac{m_1}{N}\right)\left(1-\dfrac{c_j}{N}\right)}}
\]
（調整済み残差）は，カテゴリ $j$ に属する割合を2群で比較する検定統計量
\[
z_j=\frac{\hat p_{1j}-\hat p_{2j}}{\sqrt{\hat\pi_j(1-\hat\pi_j)\left(\dfrac1{m_1}+\dfrac1{m_2}\right)}},\qquad\hat\pi_j=\frac{c_j}{N}
\]
に一致することを示せ。
\item 表1の治療群について，3つのカテゴリの調整済み残差 $r_{11},r_{12},r_{13}$ を小数第2位まで求めよ。
この結果から，上問〔1〕の検定結果にどのカテゴリが寄与しているかを述べよ。
\item 6個のセルの度数 $n_{ij}$ が互いに独立に期待値 $\mu_{ij}$ のポアソン分布に従うとし，対数線形モデル
\[
\log\mu_{ij}=\lambda+\alpha_i+\beta_j\qquad(i=1,2;\ j=1,2,3),\qquad \alpha_1=\beta_1=0
\]
を考える。対数尤度関数を書き，尤度方程式を導いて，$\mu_{ij}$ の最尤推定値が $\hat\mu_{ij}=e_{ij}$ となることを示せ。
ただし，尤度方程式の解が最尤推定値を与えることは証明なしに用いてよい。
\item 上問〔4〕のモデルの適合度を表す尤度比統計量
\[
G^2=2\sum_{i=1}^{2}\sum_{j=1}^{3}n_{ij}\log\frac{n_{ij}}{\hat\mu_{ij}}
\]
の値を表1について小数第1位まで求めよ。また，$G^2$ の漸近分布の自由度を理由とともに答え，有意水準5\%でモデルの適合度を検定せよ。
\end{enumerate}
\newpage

\qhead{3}
ある慢性疾患の治療薬の単群試験で，200人の患者について，治療開始4週後の反応 $Y_1$ と8週後の反応 $Y_2$（いずれも有効なら1，無効なら0）を調べる計画であった。
関心のある母数は8週後の有効割合 $\theta=\Prob(Y_2=1)$ である。
$Y_1$ は全員について観測されたが，一部の患者は4週後の評価の後に試験から脱落し，$Y_2$ が欠測となった。
$Y_2$ が観測されたとき $R=1$，欠測のとき $R=0$ とする。結果を表1に示す。
\begin{center}
表1：4週後の反応別の人数\par\smallskip
\begin{tabular}{crrrr}
\toprule
 & & \multicolumn{2}{c}{$Y_2$ が観測された（$R=1$）} & $Y_2$ が欠測（$R=0$）\\
\cmidrule(lr){3-4}
$Y_1$ & 人数 & 人数 & うち $Y_2=1$ & 人数\\
\midrule
1 & 80 & 72 & 63 & 8\\
0 & 120 & 60 & 18 & 60\\
\bottomrule
\end{tabular}
\end{center}
患者ごとの $(Y_1,Y_2,R)$ は独立に同一の分布に従うとし，$y=0,1$ に対して
\[
p_y=\Prob(Y_1=y),\qquad \mu_y=\Prob(Y_2=1\mid Y_1=y),\qquad
\pi(y_1,y_2)=\Prob(R=1\mid Y_1=y_1,Y_2=y_2)>0
\]
とおき，$0<p_1<1$ とする。$\pi(y_1,y_2)$ が $y_2$ によらないとき，欠測はMAR（missing at random）であるといい，このとき $\pi(y_1,y_2)=\pi_{y_1}$ と書く。
さらに $\pi_1=\pi_0$ であるとき，欠測はMCAR（missing completely at random）であるという。以下の各問に答えよ。
\begin{enumerate}
\item 表1から，次の2つの方法による $\theta$ の推定値を求めよ。
\begin{enumerate}[label=(\alph*),leftmargin=2em,itemsep=0.2em,topsep=0.2em]
\item $Y_2$ が観測された患者だけを用いる方法（完全ケース解析）。
\item $Y_2$ が欠測の患者には $Y_1$ の値を $Y_2$ の値として代入し，全員を用いる方法（LOCF）。
\end{enumerate}
\item 欠測がMARであるとする。$\Prob(Y_2=1\mid Y_1=y,R=1)=\mu_y$ であることを示せ。
また，完全ケース解析による推定量が確率収束する値 $\theta_{\mathrm{CC}}=\Prob(Y_2=1\mid R=1)$ を $p_y,\pi_y,\mu_y$ で表し，
\[
\theta_{\mathrm{CC}}-\theta=\frac{p_1p_0(\pi_1-\pi_0)(\mu_1-\mu_0)}{p_1\pi_1+p_0\pi_0}
\]
が成り立つことを示せ。この式から，完全ケース解析による推定量が $\theta$ の一致推定量になるのはどのような場合かを述べよ。
\item 欠測がMARであるとする。$\theta=p_1\mu_1+p_0\mu_0$ の右辺の各量を表1の標本割合で置き換えた推定値 $\hat\theta_{\mathrm S}$ を求めよ。
ただし，$\mu_y$ は上問〔2〕の結果に基づき，$Y_2$ が観測された患者の中での割合で置き換える。
また，$Y_1=y$ の患者のうち $Y_2$ が観測された者の割合を $\hat\pi_y$ とするとき，
\[
\hat\theta_{\mathrm{IPW}}=\frac1N\sum_{i=1}^{N}\frac{R_iY_{2i}}{\hat\pi_{Y_{1i}}}\qquad(N=200)
\]
が $\hat\theta_{\mathrm S}$ に一致することを示せ。ここで添え字 $i$ は患者を表し，$R_i=0$ の患者の項は0とする。
さらに，$\pi_y$ が既知であるとして，$\E\!\left[\dfrac{RY_2}{\pi_{Y_1}}\right]=\theta$ を示せ。
\item 欠測がMARであるとする。LOCFによる推定量が確率収束する値 $\theta_{\mathrm{LOCF}}$ を求め，$\theta_{\mathrm{LOCF}}-\theta$ を $p_y,\pi_y,\mu_y$ で表せ。
この式を用いて，表1のデータでLOCFによる推定値が上問〔3〕の $\hat\theta_{\mathrm S}$ より小さくなった理由を説明せよ。
\item MARが成り立つとは限らない場合を考える。ある定数 $k$（$0\le k\le1$）に対して
\[
\Prob(Y_2=1\mid Y_1=y,R=0)=k\,\Prob(Y_2=1\mid Y_1=y,R=1)\qquad(y=0,1)
\]
が成り立つとする。このとき表1に基づく $\theta$ の推定値を $k$ の式で表し，$k=1$ と $k=0$ がそれぞれどのような仮定に対応するかを述べよ。
また，$k$ の値を表1のデータから推定することができるか，理由とともに答えよ。
\end{enumerate}
\newpage

\qhead{4}
ある疾患の再発予防薬の無作為化比較試験で，対照群100人と治療群100人を1年間追跡したところ，再発（以下，イベントという）は対照群で40人，治療群で25人に発生した。
対照群，治療群のイベント発生割合の真値をそれぞれ $p_0,p_1$ とし，各群のイベント発生数は互いに独立に二項分布に従うとする。
治療必要数（NNT）は，リスク差 $p_0-p_1$（$>0$）の逆数 $1/(p_0-p_1)$ で定義される。以下の各問に答えよ。
\begin{enumerate}
\item リスク差 $p_0-p_1$，リスク比 $p_1/p_0$，オッズ比 $\dfrac{p_1/(1-p_1)}{p_0/(1-p_0)}$，およびNNTの推定値を求めよ。NNTは小数第2位まで求めること。
\item リスク差の標準誤差を求め，正規近似によるリスク差の95\%信頼区間を求めよ。
また，この区間の両端の逆数をとることでNNTの95\%信頼区間を求めよ。
さらに，リスク差の信頼区間が0を含む場合には，同じ方法で得られるNNTの範囲がどのようなものになるかを述べよ。
\item 各群の標本割合を $\hat p_0,\hat p_1$，人数を $n_0,n_1$ とする。デルタ法を用いて対数リスク比の推定量 $\log(\hat p_1/\hat p_0)$ の分散の近似式を導け。
これを用いてリスク比の95\%信頼区間を求めよ。
\item この疾患の患者の母集団で，治療の有無 $x$（治療あり1，なし0）とイベントの有無 $D$（発生1，非発生0）の関係がロジスティックモデル
\[
\log\frac{\Prob(D=1\mid x)}{\Prob(D=0\mid x)}=\alpha+\beta x
\]
に従うとする。この母集団から，イベントが発生した人をそれぞれ確率 $f_1$ で，発生しなかった人をそれぞれ確率 $f_0$ で，$x$ の値によらず独立に抽出するケース・コントロール研究を行う（$f_1>0,\ f_0>0$）。
抽出されたことを $S=1$ で表すとき，$\Prob(D=1\mid x,S=1)$ を求め，抽出されたデータにおいても傾きが $\beta$ のロジスティックモデルが成り立つことを示せ。切片はどのように変わるか。
この結果から，$f_1,f_0$ が未知のケース・コントロール研究のデータによって，上問〔1〕の4つの指標のうちどれが推定できるかを述べよ。
\item 治療の効果がオッズ比 $\psi=0.5$ で表され，この値は対照群のイベント発生割合 $p_0$ が異なる集団でも共通であるとする。
このとき，NNTを $p_0$ の式で表し，$p_0=0.4$ と $p_0=0.05$ のそれぞれの場合のNNTの値を求めよ。
また，NNTが最小となる $p_0$（$0<p_0<1$）の値と，そのときのNNTの値を求めよ。
これらの結果から，ある試験で得られたNNTを別の集団に当てはめるときの注意点を述べよ。
\end{enumerate}
\newpage

\qhead{5（共通問題形式）}
ある感染症の検査では，複数の人から採取した検体を混ぜ合わせた検体（以下，プール検体という）を1回の検査で調べることができる。
プール検体の検査結果は，感染者の検体が1つでも含まれていれば陽性，1つも含まれていなければ陰性となり，検査結果に誤りはないものとする。
対象とする集団では，各人は互いに独立に確率 $p$（$0<p<1$）で感染しているとし，$q=1-p$ とおく。以下の各問に答えよ。
\begin{enumerate}
\item 次の2段階の検査方式を考える。まず $k$ 人（$k\ge2$）からなるグループのプール検体を1回検査する。
結果が陰性であれば $k$ 人全員を非感染と判定して終了し，陽性であれば $k$ 人全員を1人ずつ個別に検査する。
1つのグループに必要な検査回数の期待値を求め，1人当たりの期待検査回数が
\[
h(k)=1+\frac1k-q^k
\]
となることを示せ。
\item $p=0.01$，$k=10$ のとき，$h(10)$ の値を小数第3位まで求めよ。
また，$p$ が小さいときに $q^k$ を $1-kp$ で近似した式を用いて，$h(k)$ を最小にする $k$ とそのときの $h(k)$ の値を $p$ で表せ。ただし $k$ は正の実数として扱ってよい。
\end{enumerate}
次に，感染者の割合 $p$ を推定することを目的として，$k$ 人ずつからなる $n$ 個のグループのプール検体をそれぞれ1回だけ検査し，個別の検査は行わないとする。
陽性となったプール検体の数を $X$ とする。
\begin{enumerate}[resume]
\item $X$ の従う分布を答えよ。また，$p$ の最尤推定量 $\hat p$ を求め，デルタ法を用いて $\hat p$ の分散の近似式を求めよ。
\item $k=10$，$n=50$ で，陽性となったプール検体が5個であった。$p$ の最尤推定値を小数第4位まで求めよ。
また，その標準誤差を求め，正規近似による $p$ の95\%信頼区間を求めよ（いずれも小数第4位まで）。
\item 上問〔3〕で求めた $\hat p$ の分散の近似式について，次の (a)〜(c) に答えよ。
\begin{enumerate}[label=(\alph*),leftmargin=2em,itemsep=0.3em,topsep=0.3em]
\item 同じ検査回数 $n$ で，$n$ 人を1人ずつ個別に検査して $p$ を推定する場合の推定量の分散に対する，$\hat p$ の分散の比を求めよ。
また，$p$ が小さいときにこの比がおよそ $1/k$ になることを示せ。
\item $p$ と $n$ を固定し，$k$ を正の実数として扱う。$u=-k\log q$ とおくと，$\hat p$ の分散は $(e^u-1)/u^2$ に比例することを示せ。
さらに，これを最小にする $u$ は方程式 $(u-2)e^u+2=0$ のただ1つの正の解であり，その解が $1.5<u<1.6$ を満たすことを示せ。
\item (b) の結果を踏まえて，$k$ を大きくしすぎると $p$ の推定にどのような問題が生じるかを述べよ。
\end{enumerate}
\end{enumerate}
\vfill
\begin{center}（問題は以上）\end{center}
\end{document}
